% Part: propositional-logic
% Chapter: propositional-logic 
% Section: completeness

\documentclass[../../../include/open-logic-section]{subfiles}

\begin{document}

\olfileid{pl}{prp}{com}

\olsection{પ્રતિજ્ઞાપન તર્કશાસ્ત્રની પૂર્ણતા}

\begin{defn}
\ollabel{def:MaxCon}
!!{formula}નો ગણ~$\Gamma$ \emph{મહત્તમ સુસંગત} કહેવાય,
જો તે સુસંગત હોય અને કોઈ સુસંગત ગણ~$\Delta$ માટે
$\Gamma \subseteq \Delta$ હોય ત્યારે $\Gamma = \Delta$ થાય.
\end{defn}

\begin{lem}[સત્ય સહાયક પ્રમેય]
  \ollabel{lem:truth}
$\Gamma$ મહત્તમ સુસંગત હોય, તો:
\begin{enumerate}
\item \ollabel{lem:truth-1} $!A \in \Gamma$ જો અને ફક્ત જો
  $\lnot !A \notin \Gamma$;
\item \ollabel{lem:truth-2} $\Gamma \Proves !A$ જો અને ફક્ત જો
  $!A \in \Gamma$;
\item \ollabel{lem:truth-3} $!A \lif !B \in \Gamma$ જો અને ફક્ત જો
  કાં તો $!A \notin \Gamma$ અથવા $!B \in \Gamma$.
\end{enumerate}
\end{lem}

\begin{proof}
\begin{enumerate}
\item ધારો કે $!A \in \Gamma$ હોય. જો $\lnot !A \in \Gamma$
  પણ હોય, તો $\Gamma$ અસુસંગત થાય.
  બીજી શક્યતા એ છે કે $!A$ કે $\lnot!A$ બેમાંથી એકેય
  $\Gamma$માં ન હોય. ત્યારે
  \olref[pl][axd][prv]{prop:provability-exhaustive} પ્રમાણે
  $\Gamma\cup\{!A\}$ અને $\Gamma \cup \{\lnot !A\}$માંથી
  એક સુસંગત હોય, એટલે $\Gamma$ મહત્તમ ન હોય.
  આ બંને અવલોકનો ઊલટી દિશા પણ આપે છે.
\item જો $!A \in \Gamma$ હોય, તો દેખીતી રીતે $\Gamma \Proves !A$.
  ઊલટી તરફ, ધારો કે $\Gamma \Proves !A$ હોય.
  જો $!A \notin \Gamma$ હોય, તો \olref{lem:truth-1}
  પ્રમાણે $\lnot !A \in \Gamma$ થાય.
  આમ $\Gamma \Proves !A$ અને $\Gamma \Proves \lnot !A$,
  જે સુસંગતતા સાથે વિરોધાભાસ આપે છે.
\item કસરત.
\end{enumerate}
\sourcecorrection{OLMD-344}{મૂળના બે જૂના સંદર્ભ એકેય વિસ્તરણ સુસંગત ન હોય તો મૂળ ગણ અસુસંગત થાય તે પરિણામને નિર્દેશે છે, પરંતુ આપેલા નામવાળો ઉલ્લેખ સ્થિર સ્રોતમાં મળતો નથી. તેમની જગ્યાએ હાલની સ્વયંસિદ્ધ-નિગમન રચનાના અસુસંગતતાનાં બંને વિસ્તરણ અંગેના પરિણામનો સ્પષ્ટ સંદર્ભ આપ્યો છે. આગળના નિષ્પન્નતા-અસુસંગતતા સંદર્ભને પણ તેના સાચા વિભાગ સુધી વિસ્તૃત કર્યો છે.}
\end{proof}

\olref{lem:truth}\olref{lem:truth-2}ની ડાબેથી જમણી દિશા
બતાવે છે કે $\Gamma$ \emph{નિગમન હેઠળ બંધ} છે:
દરેક $!A$ માટે, જો $\Gamma \Proves !A$ હોય, તો $!A
\in \Gamma$ થાય.

\begin{prob}
  \olref{lem:truth}ની સાબિતી પૂર્ણ કરો.
\end{prob}

\begin{thm}[પૂર્ણતા]
\ollabel{thm:completeness}
જો $\Gamma$ સુસંગત હોય, તો $\Gamma$ સંતોષ્ય છે.
\end{thm}

\begin{proof}
ભાષાનાં બધાં !!{formula}ની સંપૂર્ણ સૂચિ $!A_0$, $!A_1$,
\dots લો. !!{formula}ના ગણોનો વધતો ક્રમ
$\Gamma_0$, $\Gamma_1$, \dots નીચે મુજબ પુનરાવર્તનથી વ્યાખ્યાયિત કરીએ:
\begin{align*}
\Gamma_0 & = \Gamma\\
\Gamma_{n+1} &  =
\begin{cases}
  \Gamma_n \cup \{!A_n\} & \text{જો $\Gamma_n \cup \{!A_n\}$ સુસંગત હોય;}\\
  \Gamma_n \cup \{\lnot !A_n\} & \text{અન્યથા}.
\end{cases}
\end{align*}
પછી આ વ્યાખ્યા આપીએ:
\[
\Gamma^* = \bigcup_{0\le n}\Gamma_n.
\]
હવે સાબિતીને આગળ વધારવા ક્રમશઃ નીચેની હકીકતો સ્થાપિત કરવાની છે:
\begin{enumerate}
\item દરેક $n$ માટે ગણ~$\Gamma_n$ સુસંગત છે:
  \olref[pl][axd][prv]{prop:provability-exhaustive} વાપરી
  $n$ પર આગમન કરો;
\item $\Gamma^*$ સુસંગત છે;
\item $\Gamma^*$ મહત્તમ છે.
\end{enumerate}
પછી સત્યમૂલ્ય-નિયુક્તિ~$v$ એવી વ્યાખ્યાયિત કરીએ કે
$v(p_i) = \True$ જો અને ફક્ત જો $p_i \in \Gamma^*$.
$!A$ પર આગમનથી બતાવવાનું છે કે વધુ જટિલ !!{sentence}
માટે પણ $\Gamma^*$નું સભ્યપદ અને $v$ અનુસારનું સત્ય એક જ છે:
$\overline{v}(!A) = \True$ જો અને ફક્ત જો $!A \in \Gamma^*$.
ખાસ કરીને $v$ એ $\Gamma^*$ને સંતોષે છે;
$\Gamma \subseteq \Gamma^*$ હોવાથી $\Gamma$ પણ સંતોષ્ય છે,
જે જોઈએ છે.
\end{proof}

\begin{cor}
જો $\Gamma \Entails !A$ હોય, તો $\Gamma \Proves !A$.
\end{cor}

\begin{proof}
જો $\Gamma\Proves/ !A$ હોય, તો
\olref[pl][axd][prv]{prop:prov-incons} પ્રમાણે
$\Gamma \cup \{\lnot !A\}$ સુસંગત છે.
પૂર્ણતા પ્રમેયથી $\Gamma \cup
\{\lnot !A\}$ સંતોષ્ય છે, એટલે $\Gamma \Entails/ !A$.
\end{proof}

\begin{prop}[સઘનતા પ્રમેય]
\ollabel{prop:compactness}
$\Gamma$ સંતોષ્ય છે જો અને ફક્ત જો $\Gamma$નો દરેક
\emph{સાંત} ઉપગણ~$\Gamma_0$ સંતોષ્ય હોય.
\end{prop}

\begin{proof}
$\Gamma$ અસંતોષ્ય છે જો અને ફક્ત જો તે અસુસંગત હોય;
જો અને ફક્ત જો $\Gamma$નો કોઈ સાંત ઉપગણ~$\Gamma_0$
અસુસંગત હોય; જો અને ફક્ત જો $\Gamma$નો કોઈ સાંત
ઉપગણ~$\Gamma_0$ અસંતોષ્ય હોય.
\end{proof}

\begin{cor}
$\Gamma \Entails !A$ જો અને ફક્ત જો $\Gamma$ના
કોઈ સાંત ઉપગણ~$\Gamma_0$ માટે $\Gamma_0 \Entails !A$.
\end{cor}

\end{document}

